\documentclass[10pt,a4paper]{amsart}
\usepackage{amsmath,amssymb,amsthm,graphicx}
\usepackage{verbatim}
\newenvironment{reference}{\comment}{\endcomment}
%\newenvironment{reference}{}{}
\newenvironment{slogan}{\comment}{\endcomment}
\newenvironment{history}{\comment}{\endcomment}

% For commutative diagrams we use Xy-pic
\usepackage[all]{xy}

% We use 2cell for 2-commutative diagrams.
\xyoption{2cell}
% Theorem environments.
%
\theoremstyle{plain}
\newtheorem{theorem}[subsection]{Theorem}
\newtheorem{proposition}[subsection]{Proposition}
\newtheorem{lemma}[subsection]{Lemma}

\theoremstyle{definition}
\newtheorem{definition}[subsection]{Definition}
\newtheorem{example}[subsection]{Example}
\newtheorem{exercise}[subsection]{Exercise}
\newtheorem{situation}[subsection]{Situation}

\theoremstyle{remark}
\newtheorem{remark}[subsection]{Remark}
\newtheorem{remarks}[subsection]{Remarks}

\numberwithin{equation}{subsection}

% Macros
%
\def\lim{\mathop{\mathrm{lim}}\nolimits}
\def\colim{\mathop{\mathrm{colim}}\nolimits}
\def\Spec{\mathop{\mathrm{Spec}}}
\def\Hom{\mathop{\mathrm{Hom}}\nolimits}
\def\Ext{\mathop{\mathrm{Ext}}\nolimits}
\def\SheafHom{\mathop{\mathcal{H}\!\mathit{om}}\nolimits}
\def\SheafExt{\mathop{\mathcal{E}\!\mathit{xt}}\nolimits}
\def\Sch{\mathit{Sch}}
\def\Mor{\mathop{\mathrm{Mor}}\nolimits}
\def\Ob{\mathop{\mathrm{Ob}}\nolimits}
\def\Sh{\mathop{\mathit{Sh}}\nolimits}
\def\NL{\mathop{N\!L}\nolimits}
\def\CH{\mathop{\mathrm{CH}}\nolimits}
\def\proetale{{pro\text{-}\acute{e}tale}}
\def\etale{{\acute{e}tale}}
\def\QCoh{\mathit{QCoh}}
\def\Ker{\mathop{\mathrm{Ker}}}
\def\Im{\mathop{\mathrm{Im}}}
\def\Coker{\mathop{\mathrm{Coker}}}
\def\Coim{\mathop{\mathrm{Coim}}}

% Boxtimes
%
\DeclareMathSymbol{\boxtimes}{\mathbin}{AMSa}{"02}

%
% Macros for moduli stacks/spaces
%
\def\QCohstack{\mathcal{QC}\!\mathit{oh}}
\def\Cohstack{\mathcal{C}\!\mathit{oh}}
\def\Spacesstack{\mathcal{S}\!\mathit{paces}}
\def\Quotfunctor{\mathrm{Quot}}
\def\Hilbfunctor{\mathrm{Hilb}}
\def\Curvesstack{\mathcal{C}\!\mathit{urves}}
\def\Polarizedstack{\mathcal{P}\!\mathit{olarized}}
\def\Complexesstack{\mathcal{C}\!\mathit{omplexes}}
% \Pic is the operator that assigns to X its picard group, usage \Pic(X)
% \Picardstack_{X/B} denotes the Picard stack of X over B
% \Picardfunctor_{X/B} denotes the Picard functor of X over B
\def\Pic{\mathop{\mathrm{Pic}}\nolimits}
\def\Picardstack{\mathcal{P}\!\mathit{ic}}
\def\Picardfunctor{\mathrm{Pic}}
\def\Deformationcategory{\mathcal{D}\!\mathit{ef}}
\usepackage{cleveref}
\crefname{originalreference}{original reference}{original references}
\Crefname{originalreference}{Original reference}{Original references}
\makeatletter
\newlabel{lemma-geometric-lift-to-usual}{{05VN}{source}{}{sourcepointer-1}{}}
\newlabel{lemma-geometric-lift-to-usual@cref}{{[originalreference][0][]05VN}{[1][0][]source}{}{}{}}
\makeatother
\crefname{originalreference}{original reference}{original references}
\Crefname{originalreference}{Original reference}{Original references}
\makeatletter
\newlabel{lemma-stalk-pullback}{{008H}{source}{}{sourcepointer-2}{}}
\newlabel{lemma-stalk-pullback@cref}{{[originalreference][0][]008H}{[1][0][]source}{}{}{}}
\makeatother
\crefname{originalreference}{original reference}{original references}
\Crefname{originalreference}{Original reference}{Original references}
\makeatletter
\newlabel{lemma-stalk-quasi-coherent}{{05VP}{source}{}{sourcepointer-3}{}}
\newlabel{lemma-stalk-quasi-coherent@cref}{{[originalreference][0][]05VP}{[1][0][]source}{}{}{}}
\makeatother
\crefname{originallemma}{original lemma}{original lemmas}
\Crefname{originallemma}{Original lemma}{Original lemmas}
\makeatletter
\newlabel{sheaves-lemma-stalk-pullback-modules}{{0098}{source}{}{sourcepointer-4}{}}
\newlabel{sheaves-lemma-stalk-pullback-modules@cref}{{[originallemma][0][]0098}{[1][0][]source}{}{}{}}
\makeatother
\title{Stalks of the pullback of a quasi-coherent module at a geometric point}
\author{The Stacks Project Authors}
\date{stacks-project:spaces-properties.tex:local-snapshot-20261012; GNU FDL 1.2 or later}
\begin{document}
\maketitle
\noindent\textit{Editorial note.} All delivered mathematical text is the author's own, copied from the pinned Stacks Project chapter file spaces-properties.tex; no statement or proof step was rewritten, reordered or removed.\par\medskip
\noindent\textit{Editorial note.} The counted proof is the author's complete argument: the reduction to the tensor product description of the pullback together with the commutation of stalks with tensor products, and the explicit computation on a compatible square of presentations with a lifted geometric point.\par\medskip
\noindent\textit{Editorial note.} The label lemma-stalk-pullback-quasi-coherent names the counted statement, and the proof environment immediately following it in the source is delivered verbatim as its proof.\par\medskip
\noindent\textit{Editorial note.} The counted statement is carried by the author's own lemma environment in the chapter source.\par\medskip
\noindent\textit{Editorial note.} Pre-open gate: python operations/agent-authoring/tools/precard\_gate.py --source spaces-properties.tex --label lemma-stalk-pullback-quasi-coherent --start 5547 --end 5606 --local-id 28050 returned RESULT=OK (L28050 free, occupancy FREE, public interval clear). Candidate taken from sources/stacks-spaces-properties/candidates-halfA.json.\par\medskip
\noindent\textit{Editorial note.} The statement and proof are copied from the pinned chapter file; the cited inputs are kept as external pointers with their canonical Stacks tags.\par\medskip
\noindent\textit{Allowed standard prerequisites (editorial).}\begin{itemize}
\item Description of the pullback of a module as a tensor product. Same chapter, the definition of the pullback of a module as the tensor product of the pullback of sheaves with the structure sheaf of the source. Supplies the reduction of the statement to the tensor product description.
\item Stalks of pullbacks and of quasi-coherent modules. Same chapter, the lemma describing the stalks of the pullback of a sheaf of sets, and the lemma identifying the stalk of a quasi-coherent module at a geometric point with the tensor product of the stalk at the underlying point with the local ring. Supplies the two identifications used in the displayed computation.
\item Lifting a geometric point along a surjective etale morphism. Same chapter, the lemma allowing a geometric point of an algebraic space to be lifted to a scheme presenting it. Supplies the geometric point u of U mapping to the given geometric point of X.
\end{itemize}\medskip
\ifdefined\hypertarget\hypertarget{sourcepointer-1}{}\fi
\noindent\textit{Source reference.} Original reference 05VN: sources/stacks-spaces-properties/pinned-original-source.tex, source lines 3381--3381. Prerequisite pointer invoked at proof line(s) 5579; it names a standard fact of the same chapter that the counted argument uses, and every step of the counted proof built on it is retained verbatim in the delivered body.\par\medskip
\ifdefined\hypertarget\hypertarget{sourcepointer-2}{}\fi
\noindent\textit{Source reference.} Original reference 008H: sources/stacks-spaces-properties/pinned-original-source.tex, source lines 3594--3594. Prerequisite pointer invoked at proof line(s) 5567; it names a standard fact of the same chapter that the counted argument uses, and every step of the counted proof built on it is retained verbatim in the delivered body.\par\medskip
\ifdefined\hypertarget\hypertarget{sourcepointer-3}{}\fi
\noindent\textit{Source reference.} Original reference 05VP: sources/stacks-spaces-properties/pinned-original-source.tex, source lines 5505--5505. Prerequisite pointer invoked at proof line(s) 5602; it names a standard fact of the same chapter that the counted argument uses, and every step of the counted proof built on it is retained verbatim in the delivered body.\par\medskip
\ifdefined\hypertarget\hypertarget{sourcepointer-4}{}\fi
\noindent\textit{Source reference.} Original lemma 0098: sources/stacks-chow/tags.txt, source lines 344--344. Prerequisite pointer invoked at proof line(s) 5605; it names a standard fact of a sibling chapter that the counted argument uses, and every step of the counted proof built on it is retained verbatim in the delivered body.\par\medskip
\begin{lemma}
\label{lemma-stalk-pullback-quasi-coherent}
Let $S$ be a scheme. Let $f : X \to Y$ be a morphism of algebraic spaces
over $S$. Let $\mathcal{G}$ be a quasi-coherent $\mathcal{O}_Y$-module.
Let $\overline{x}$ be a geometric point of $X$ and let
$\overline{y} = f \circ \overline{x}$ be the image in $Y$.
Then there is a canonical isomorphism
$$
(f^*\mathcal{G})_{\overline{x}} =
\mathcal{G}_{\overline{y}} \otimes_{\mathcal{O}_{Y, \overline{y}}}
\mathcal{O}_{X, \overline{x}}
$$
of the stalk of the pullback with the tensor product of the stalk
with the local ring of $X$ at $\overline{x}$.
\end{lemma}

\begin{proof}
Since $f^*\mathcal{G} =
f_{small}^{-1}\mathcal{G} \otimes_{f_{small}^{-1}\mathcal{O}_Y} \mathcal{O}_X$
this follows from the description of stalks of pullbacks in
Lemma \ref{lemma-stalk-pullback}
and the fact that taking stalks commutes with tensor products.
A more direct way to see this is as follows.
Choose a commutative diagram
$$
\xymatrix{
U \ar[d]_p \ar[r]_\alpha  & V \ar[d]^q \\
X \ar[r]^a & Y
}
$$
where $U$ and $V$ are schemes, and $p$ and $q$ are surjective \'etale.
By
Lemma \ref{lemma-geometric-lift-to-usual}
we can choose a geometric point $\overline{u}$ of $U$ such that
$\overline{x} = p \circ \overline{u}$. Set
$\overline{v} = \alpha \circ \overline{u}$.
Then we see that
\begin{align*}
(f^*\mathcal{G})_{\overline{x}} & =
(p^*f^*\mathcal{G})_u \otimes_{\mathcal{O}_{U, u}}
\mathcal{O}_{X, \overline{x}} \\
& = (\alpha^*q^*\mathcal{G})_u \otimes_{\mathcal{O}_{U, u}}
\mathcal{O}_{X, \overline{x}} \\
& = (q^*\mathcal{G})_v \otimes_{\mathcal{O}_{V, v}}
\mathcal{O}_{U, u} \otimes_{\mathcal{O}_{U, u}}
\mathcal{O}_{X, \overline{x}} \\
& = (q^*\mathcal{G})_v \otimes_{\mathcal{O}_{V, v}}
\mathcal{O}_{X, \overline{x}} \\
& = (q^*\mathcal{G})_v \otimes_{\mathcal{O}_{V, v}}
\mathcal{O}_{Y, \overline{y}} \otimes_{\mathcal{O}_{Y, \overline{y}}}
\mathcal{O}_{X, \overline{x}} \\
& = \mathcal{G}_{\overline{y}} \otimes_{\mathcal{O}_{Y, \overline{y}}}
\mathcal{O}_{X, \overline{x}}
\end{align*}
Here we have used
Lemma \ref{lemma-stalk-quasi-coherent} (twice)
and the corresponding result for pullbacks of quasi-coherent sheaves
on schemes, see
Sheaves, Lemma \ref{sheaves-lemma-stalk-pullback-modules}.
\end{proof}

\end{document}
